\subsection{From Polarization to Qubits}

The example of photon polarization (\autopageref{ex:polarization}) shows that a physical system can be described using two basic states and their linear combinations. We now \textbf{abstract} this idea and introduce a \textbf{general model that does not depend on a specific physical system}.

\highspace
\begin{flushleft}
    \textcolor{Green3}{\faIcon{question-circle} \textbf{From Physical States to Abstract States}}
\end{flushleft}
In the polarization example, we used:
\begin{equation*}
    \rightarrow \quad \text{and} \quad \uparrow
\end{equation*}
As the two basic states of the system. We now replace them with an abstract notation:
\begin{equation*}
    \rightarrow \;\rightsquigarrow\; \ket{0}
    \hspace{2em}
    \uparrow \;\rightsquigarrow\; \ket{1}
\end{equation*}
These are called the \textbf{computational basis states}. They are just symbols (used in Dirac notation) that represent the two basic states of the system, without referring to any specific physical implementation.

\highspace
\begin{flushleft}
    \textcolor{Green3}{\faIcon{book} \textbf{Definition of a Qubit}}
\end{flushleft}
A \definition{Qubit} is a \textbf{two-level quantum system} whose state can be written as:
\begin{equation*}
    \ket{\psi} = a \ket{0} + b \ket{1}
\end{equation*}
It is the \textbf{unit of quantum information}, analogous to a classical bit. However, unlike a classical bit that can only be in one of two states (0 or 1), a qubit can exist in \textbf{infinitely many possible states}. Indeed, any normalized vector of the form:
\begin{equation*}
    \ket{\psi} = a \ket{0} + b \ket{1}
\end{equation*}
Represents a valid state of a qubit. Some \textbf{properties} of the qubit state are:
\begin{itemize}
    \item The coefficients $a$ and $b$ are \textbf{complex numbers} (probability amplitudes). They determine the outcome of a measurement:
    \begin{itemize}
        \item $\left|a\right|^2$ is the probability of measuring the state $\ket{0}$.
        \item $\left|b\right|^2$ is the probability of measuring the state $\ket{1}$.
    \end{itemize}
    Therefore, even though a qubit can be in infinitely many possible states, a measurement (in this basis) can only produce \textbf{two possible outcomes}: $\ket{0}$ or $\ket{1}$.
    
    \item The state is \textbf{normalized}:
    \begin{equation*}
        \left|a\right|^2 + \left|b\right|^2 = 1
    \end{equation*}
    Because the probabilities must sum to 1 (i.e., we must get either $\ket{0}$ or $\ket{1}$ when we measure).

    \item The states $\ket{0}$ and $\ket{1}$ form a \textbf{basis} of the vector space describing the qubit.

    \highspace
    \textcolor{Green3}{\faIcon{question-circle} \textbf{What is a \emph{basis}?}} A \definition{Basis} is a set of vectors that allows us to \textbf{represent any vector in the space as a linear combination of them}. In quantum computing, $\ket{0}$ and $\ket{1}$ are the \textbf{basis vectors} and any qubit state is built from them.

    \highspace
    In particular, the states $\ket{0}$ and $\ket{1}$ form an \definition{Orthonormal Basis}, meaning that:
    \begin{itemize}
        \item They are \textbf{orthogonal}: $\braket{0}{1} = 0$.
        
        \textcolor{Green3}{\faIcon{question-circle} \textbf{Why is orthogonality important?}} Orthogonality means that the two basis states are completely distinguishable by a measurement (not overlapping). If we measure a qubit in the state $\ket{0}$, we will never get $\ket{1}$, and vice versa.


        \item They are \textbf{normalized}: $\braket{0}{0} = \braket{1}{1} = 1$.
        
        \textcolor{Green3}{\faIcon{question-circle} \textbf{Why is normalization important?}} Normalization ensures that the probabilities of measurement outcomes are well-defined and sum to 1. It also means that the basis states have unit length in the vector space, which is a fundamental property for quantum states.
    \end{itemize}
\end{itemize}

\highspace
\begin{flushleft}
    \textcolor{Green3}{\faIcon{question-circle} \textbf{Why This Abstraction Matters}}
\end{flushleft}
This abstraction is fundamental because it allows us to:
\begin{itemize}
    \item Describe \textbf{any} two-level quantum system (not only photons);
    \item Use a unified mathematical framework based on vectors;
    \item Apply the tools of linear algebra to quantum systems.
\end{itemize}
From now on, we will work with qubits using Dirac notation, without referring to a specific physical implementation.

\highspace
\begin{examplebox}[: Different Choices of Basis]
    The computational basis $\{\ket{0}, \ket{1}\}$ is not the only possible choice. A qubit can be described using \textbf{any orthonormal basis}.

    \highspace
    For example, another important basis is the \textbf{Hadamard basis}, defined as:
    \begin{equation*}
        \ket{+} = \frac{1}{\sqrt{2}}
        \begin{bmatrix}
            1 \\ 1
        \end{bmatrix}
        \hspace{2em}
        \ket{-} = \frac{1}{\sqrt{2}}
        \begin{bmatrix}
            1 \\ -1
        \end{bmatrix}
    \end{equation*}
    This shows that the same quantum state can be represented in different bases, depending on the context.
\end{examplebox}

\begin{examplebox}[: Orthogonality of the Hadamard Basis]
    \hl{Show that the Hadamard basis states $\ket{+}$ and $\ket{-}$ are orthogonal.}

    \begin{proof}
        Recall:
        \begin{equation*}
            \ket{+} = \frac{1}{\sqrt{2}}
            \begin{bmatrix}
                1 \\ 1
            \end{bmatrix}
            \hspace{2em}
            \ket{-} = \frac{1}{\sqrt{2}}
            \begin{bmatrix}
                1 \\ -1
            \end{bmatrix}
        \end{equation*}
        Orthogonal means that their inner product is zero. So, we need to compute $\langle + | - \rangle$ and show that it equals zero:
        \begin{equation*}
            \langle + | - \rangle =
            \frac{1}{\sqrt{2}} \begin{bmatrix}
                1 & 1
            \end{bmatrix}
            \cdot
            \frac{1}{\sqrt{2}}
            \begin{bmatrix}
                1 \\ -1
            \end{bmatrix}
            = \dfrac{1}{2}\left(1 - 1\right) = 0
        \end{equation*}
        Therefore, $\ket{+}$ and $\ket{-}$ are orthogonal.
    \end{proof}
\end{examplebox}

\begin{examplebox}[: Change of Basis from Standard $\leftrightarrow$ Hadamard]
    Given a qubit state in the standard basis:
    \begin{equation*}
        \ket{\psi} = a \ket{0} + b \ket{1}
    \end{equation*}
    We want to express it in the Hadamard basis:
    \begin{equation*}
        \ket{\psi} = x \ket{+} + y \ket{-}
    \end{equation*}

    \begin{proof}
        Using the definitions of $\ket{+}$ and $\ket{-}$:
        \begin{equation*}
            \ket{+} = \frac{1}{\sqrt{2}}
            \begin{bmatrix}
                1 \\ 1
            \end{bmatrix}
            \quad
            \ket{-} = \frac{1}{\sqrt{2}}
            \begin{bmatrix}
                1 \\ -1
            \end{bmatrix}
        \end{equation*}
        We can write $\ket{\psi}$ in terms of $\ket{+}$ and $\ket{-}$:
        \begin{equation*}
            \ket{\psi} = a \ket{0} + b \ket{1}
            = a \cdot
            \begin{bmatrix}
                1 \\ 0
            \end{bmatrix}
            + b \cdot
            \begin{bmatrix}
                0 \\ 1
            \end{bmatrix}
            =
            \begin{bmatrix}
                a \\ b
            \end{bmatrix}
        \end{equation*}
    
        Now, we express $\ket{\psi}$ as a linear combination of $\ket{+}$ and $\ket{-}$:
        \begin{equation*}
            \ket{\psi} = x \ket{+} + y \ket{-}
        \end{equation*}

        Substituting the definitions of $\ket{+}$ and $\ket{-}$:
        \begin{equation*}
            = x \cdot \frac{1}{\sqrt{2}}
            \begin{bmatrix}
                1 \\ 1
            \end{bmatrix}
            + y \cdot \frac{1}{\sqrt{2}}
            \begin{bmatrix}
                1 \\ -1
            \end{bmatrix}
        \end{equation*}

        Combining the terms:
        \begin{equation*}
            = \frac{1}{\sqrt{2}}
            \begin{bmatrix}
                x + y \\ x - y
            \end{bmatrix}
        \end{equation*}
        We want this to equal $\begin{bmatrix}
            a \\ b
        \end{bmatrix}$, so we set:
        \begin{equation*}
            \frac{1}{\sqrt{2}}
            \begin{bmatrix}
                x + y \\ x - y
            \end{bmatrix}
            =
            \begin{bmatrix}
                a \\ b
            \end{bmatrix}
        \end{equation*}
        This gives us the system of equations:
        \begin{equation*}
            x + y = \sqrt{2} a,
            \qquad
            x - y = \sqrt{2} b
        \end{equation*}
        Adding the two equations:
        \begin{equation*}
            2x = \sqrt{2} (a + b) \implies x = \frac{a + b}{\sqrt{2}}
        \end{equation*}
        Subtracting the second from the first:
        \begin{equation*}
            2y = \sqrt{2} (a - b) \implies y = \frac{a - b}{\sqrt{2}}
        \end{equation*}
        It follows that:
        \begin{equation*}
            \ket{\psi}
            = x \ket{+} + y \ket{-}
            = \frac{a + b}{\sqrt{2}} \ket{+} + \frac{a - b}{\sqrt{2}} \ket{-}
        \end{equation*}
    \end{proof}
\end{examplebox}

\begin{examplebox}[: Orthogonal State Construction]
    Given:
    \begin{equation*}
        \ket{v} = a \ket{0} + b \ket{1}
    \end{equation*}
    Define:
    \begin{equation*}
        \ket{w} = \bar{b} \ket{0} - \bar{a} \ket{1}
    \end{equation*}
    Show that $\ket{v}$ and $\ket{w}$ are orthogonal.

    \begin{proof}
        To show that $\ket{v}$ and $\ket{w}$ are orthogonal, we need to compute their inner product $\braket{w}{v}$ and show that it equals zero. First, we write the inner product converting $\ket{w}$ to its corresponding bra:
        \begin{equation*}
            \ket{w} = \bar{b} \ket{0} - \bar{a} \ket{1} \implies \bra{w} = b \bra{0} - a \bra{1}
        \end{equation*}
        Now we compute:
        \begin{equation*}
            \braket{w}{v}
            =
            \left(b \bra{0} - a \bra{1}\right) \cdot \left(a \ket{0} + b \ket{1}\right)
        \end{equation*}
        Expanding this expression:
        \begin{equation*}
            = b a \braket{0}{0}
            + b b \braket{0}{1}
            - a a \braket{1}{0}
            - a b \braket{1}{1}
        \end{equation*}
        Using the orthonormality of the basis states:
        \begin{equation*}
            \braket{0}{0} = \braket{1}{1} = 1,
            \quad
            \braket{0}{1} = \braket{1}{0} = 0
        \end{equation*}
        We get:
        \begin{equation*}
            \braket{w}{v} = b a \cancelto{1}{\braket{0}{0}}
            + \cancel{b b \braket{0}{1}}
            - \cancel{a a \braket{1}{0}}
            - a b \cancelto{1}{\braket{1}{1}}
            =
            b a - a b
        \end{equation*}
        Since $a$ and $b$ are complex numbers, multiplication is commutative, so $ba = ab$. Therefore:
        \begin{equation*}
            \braket{w}{v} = ba - ab = 0
        \end{equation*}
        Hence, $\ket{v}$ and $\ket{w}$ are orthogonal.
    \end{proof}
\end{examplebox}